\documentclass[10pt,a4paper]{article}
\usepackage[utf8]{inputenc}
\usepackage[a4paper,left=3cm,right=2cm]{geometry}
\usepackage{amsmath}
\usepackage[thmmarks, amsmath, thref]{ntheorem}
\usepackage{amsfonts}
\usepackage{amssymb}
\usepackage{fancyhdr}
\pagestyle{fancy}
\usepackage{anysize}
\usepackage{graphicx}
\usepackage{float}
\usepackage{hyperref}
\usepackage{multirow}
\usepackage{enumitem}
\newtheorem{theorem}{Theorem}
\theoremsymbol{\ensuremath{\blacksquare}}
\newtheorem*{solution}{Solución:}
\usepackage{listings}
\usepackage{xcolor}

% Margenes y formalidades

\textwidth 6.25in % Margen Derecho y ancho del texto
\textheight 8.5in
\oddsidemargin 0in % Margen Izquierdo
\headheight 0in % Largo del texto
\leftmargin 1in
\parindent 0em % Salto entre la indentación (Sangría)https://www.overleaf.com/4519127669sdpptkhvsbfn
\topmargin -0.5in % Margen Superior
\parskip 0.5ex % Salto entre parrafos
\baselineskip 1pt

%\lhead{\includegraphics[scale=0.2]{logo.png}} % texto izquierda de la cabecera
%\chead{}                                      % texto centro de la cabecera
%\rhead{\includegraphics[scale=0.2]{FIG/logo-mba-2016.png}} % número de página a la derecha

\renewcommand{\headrulewidth}{0.4pt} % grosor de la línea de la cabecera
\renewcommand{\footrulewidth}{0.4pt} % grosor de la línea del pie

\begin{document}



\pagebreak
\vspace*{4mm} 

\begin{center}
\begin{huge}
\textbf{\\Ejercicios Econometría 2}
\end{Large}\\
%{\large 2023-2}\\
\end{center}

\begin{enumerate}

  \item  Se cuenta con la siguiente información para 5 individuos, donde $Y$ es la variable dependiente y $X_1$, $X_2$ son variables explicativas:

\[
\begin{array}{c|ccc}
\hline
\text{Individuo} & X_1 & X_2 & Y \\
\hline
1 & 4 & 1 & 5 \\
2 & 3 & 2 & 7 \\
3 & 5 & 3 & 10 \\
4 & 7 & 4 & 13 \\
5 & 6 & 4 & 12 \\
\hline
\end{array}
\]

\begin{enumerate}
  \item \textbf{(Modelo de Regresión Lineal Simple)} Estime el modelo $Y = \beta_0 + \beta_1\,X_1 + \varepsilon$, ignorando la variable $X_2$. Calcule los coeficientes, el $R^2$, el $R^2$ ajustado y las sumas de cuadrados (SST, SSR y SSE).



\textbf{Soluci\'on}


Deseamos estimar 
\[
Y = \beta_0 + \beta_1 X_1 + \varepsilon,
\]
ignorando la variable $X_2$.

\paragraph{1.1 Media de $X_1$ y de $Y$.}

\[
\bar{X}_1
= \frac{4+3+5+7+6}{5}
= \frac{25}{5}
= 5,
\quad
\bar{Y}
= \frac{5+7+10+13+12}{5}
= \frac{47}{5}
= 9.4.
\]

\paragraph{1.2 Varianza, covarianza y c\'alculo de $\beta_1$ y $\beta_0$.}

Para la regresión simple:
\[
\beta_1
= \frac{\mathrm{Cov}(X_1,Y)}{\mathrm{Var}(X_1)},
\quad
\beta_0
= \bar{Y} - \beta_1\,\bar{X}_1.
\]

Construyamos la tabla con $(X_1-\bar{X}_1)$, $(Y-\bar{Y})$, y sus productos.

\[
\begin{array}{|c|c|c|c|c|c|c|c|}
\hline
\text{Ind} & X_1 & Y 
& X_1-\bar{X}_1 & Y-\bar{Y}
& (X_1 - \bar{X}_1)^2 & (Y - \bar{Y})^2
& (X_1-\bar{X}_1)(Y-\bar{Y})\\
\hline
1 & 4 & 5 
& 4-5=-1
& 5-9.4=-4.4
& 1
& 19.36
& (+4.4) \\[6pt]
2 & 3 & 7
& 3-5=-2
& 7-9.4=-2.4
& 4
& 5.76
& (+4.8) \\[6pt]
3 & 5 & 10
& 5-5=0
& 10-9.4=0.6
& 0
& 0.36
& 0 \\[6pt]
4 & 7 & 13
& 7-5=2
& 13-9.4=3.6
& 4
& 12.96
& 7.2 \\[6pt]
5 & 6 & 12
& 6-5=1
& 12-9.4=2.6
& 1
& 6.76
& 2.6 \\[6pt]
\hline
\text{Suma} & & & & 
& 10
& 45.20
& 19.0 \\[6pt]
\hline
\end{array}
\]

De la última fila:
\[
\sum (X_1-\bar{X}_1)^2 = 10,
\quad
\sum (Y - \bar{Y})^2 = 45.20,
\quad
\sum (X_1-\bar{X}_1)(Y-\bar{Y}) = 19.0.
\]
Entonces
\[
\mathrm{Var}(X_1)
= \frac{1}{n-1}\sum (X_1-\bar{X}_1)^2
= \frac{10}{4}
= 2.5,
\quad
\mathrm{Cov}(X_1,Y)
= \frac{19.0}{4}
= 4.75.
\]
\[
\beta_1
= \frac{4.75}{2.5}
= 1.90,
\quad
\beta_0
= 9.4 - (1.90)(5)
= 9.4 - 9.5
= -0.1.
\]

\[
\boxed{
\widehat{\beta}_0 \approx -0.10,
\quad
\widehat{\beta}_1 \approx 1.90.
}
\]


\[
\text{SST}
= \sum (Y_i - \bar{Y})^2
= 45.20.
\]
El valor ajustado: 
\[
\hat{Y}_i
= -0.1 +1.90\,X_{1i}.
\]
Calculamos los residuos $e_i = Y_i - \hat{Y}_i$:

\[
\begin{array}{|c|c|c|c|c|c|}
\hline
\text{Ind} & X_1 & Y 
& \hat{Y}_i = -0.1 +1.90X_1
& e_i= Y-\hat{Y}_i
& e_i^2
\\\hline
1 & 4 & 5 
& -0.1 +1.90\cdot4 =-0.1 +7.6=7.5
& 5-7.5=-2.5
& 6.25
\\
2 & 3 &7
& -0.1 +1.90\cdot3 =-0.1 +5.7=5.6
& 7-5.6=1.4
& 1.96
\\
3 & 5 &10
& -0.1 +1.90\cdot5 =-0.1 +9.5=9.4
& 10-9.4=0.6
& 0.36
\\
4 & 7 &13
& -0.1 +1.90\cdot7 =-0.1 +13.3=13.2
& 13-13.2=-0.2
& 0.04
\\
5 & 6 &12
& -0.1 +1.90\cdot6 =-0.1 +11.4=11.3
& 12-11.3=0.7
& 0.49
\\
\hline
& & & & \sum e_i^2= & 9.10
\\\hline
\end{array}
\]
\[
\text{SSR} = \sum e_i^2 \approx 9.10,
\quad
\text{SSE} = \text{SST} - \text{SSR} = 45.20 -9.10=36.10.
\]
\[
R^2
= 1 - \frac{\text{SSR}}{\text{SST}}
= 1 - \frac{9.10}{45.20}
\approx 0.7991.
\]
Es decir, el \textbf{modelo simple} explica alrededor de $80\%$ de la variabilidad de $Y$.

\paragraph{1.3.1 $R^2$ ajustado (modelo simple).}

Con $n=5$, $k=1$:
\[
R_{\text{ajustado}}^2
= 1 - \frac{\text{SSR}/(n - k -1)}{\text{SST}/(n-1)}
= 1 - \frac{9.10/3}{45.20/4}
= 1 - \frac{3.0333}{11.30}
\approx 1 -0.2684
= 0.7316.
\]


\[
Y \approx -0.10 +1.90\,X_1,
\quad
R^2 \approx 0.799,
\quad
R_{\text{ajustado}}^2 \approx 0.732.
\]

  \item \textbf{Discuta la posibilidad de sesgo} al omitir $X_2$.

\\

\textbf{Solución}

La variable $X_2$ podr\'ia ser relevante. Si $X_2$ est\'a correlacionada con $X_1$ o explica parte de la variabilidad de $Y$, omitirla puede causar que la pendiente estimada $\beta_1$ en la regresión simple sea \emph{sesgada}. Para observar si este es el caso, estimamos ahora el modelo con ambas variables y vemos si $R^2$ crece y si $\beta_1$ se modifica.

  \item \textbf{(Modelo de Regresión Lineal Múltiple)} Estime el modelo $Y = \beta_0 + \beta_1 X_1 + \beta_2 X_2 + \varepsilon$, calcule el nuevo $R^2$ (y las sumas de cuadrados) y compare con el modelo simple. Hint: 
  \[
(X^TX)^{-1}
= \frac{1}{95}
\begin{bmatrix}
281 & -72 & 35\\
-72 & 34 & -35\\
35 & -35 & 50
\end{bmatrix}.
\]


\textbf{Solución}

El modelo completo:
\[
Y = \beta_0 + \beta_1\,X_1 + \beta_2\,X_2 + \varepsilon.
\]
En notaci\'on matricial:
\[
\hat{\boldsymbol{\beta}}
= (\mathbf{X}^\top \mathbf{X})^{-1} \, (\mathbf{X}^\top \mathbf{Y}),
\]
donde
\[
\mathbf{X}
=
\begin{bmatrix}
1 & X_1 & X_2
\end{bmatrix},
\quad
\mathbf{Y}
=
\begin{bmatrix}
Y
\end{bmatrix}.
\]


\[
\mathbf{X}
=
\begin{bmatrix}
1 & 4 & 1\\
1 & 3 & 2\\
1 & 5 & 3\\
1 & 7 & 4\\
1 & 6 & 4
\end{bmatrix},
\quad
\mathbf{Y}
=
\begin{bmatrix}
5\\
7\\
10\\
13\\
12
\end{bmatrix}.
\]



Sea
\[
\mathbf{X}^\top \mathbf{Y}
=
\begin{bmatrix}
47\\
254\\
149
\end{bmatrix}.
\]

Primero calculamos 
\(\mathrm{adj}(A)\cdot b\),
sin dividir todavía por 95:
\[
\begin{bmatrix}
281 & -72 & 35\\
-72 & 34 & -35\\
35 & -35 & 50
\end{bmatrix}
\begin{bmatrix}
47\\
254\\
149
\end{bmatrix}
=
\begin{bmatrix}
\text{fila1}\\
\text{fila2}\\
\text{fila3}
\end{bmatrix}.
\]
\[
\text{(fila 1)} 
= 281\cdot47 + (-72)\cdot254 + 35\cdot149.
\]
\[
281\cdot47=13207,\quad
(-72)\cdot254=-18288,\quad
35\cdot149=5215.
\]
Suma:
\[
13207 -18288 +5215= 13207 +5215=18422,\quad 18422-18288=134.
\]
\[
\text{(fila 2)}
= (-72)\cdot47 +34\cdot254 +(-35)\cdot149.
\]
\[
(-72)\cdot47=-3384,\quad
34\cdot254=8636,\quad
(-35)\cdot149=-5215.
\]
Suma:
\[
-3384+8636=5252,\quad 5252-5215=37.
\]
\[
\text{(fila 3)}
= 35\cdot47 +(-35)\cdot254 +50\cdot149.
\]
\[
35\cdot47=1645,\quad
(-35)\cdot254=-8890,\quad
50\cdot149=7450.
\]
Suma:
\[
1645-8890 +7450=9095-8890=205.
\]
El vector intermedio es
\[
\begin{bmatrix}
134\\
37\\
205
\end{bmatrix}.
\]
Dividimos por $95$:
\[
\hat{\boldsymbol{\beta}}
=
\frac{1}{95}
\begin{bmatrix}
134\\
37\\
205
\end{bmatrix}
=
\begin{bmatrix}
134/95\\
37/95\\
205/95
\end{bmatrix}
\approx
\begin{bmatrix}
1.41\\
0.39\\
2.16
\end{bmatrix}.
\]
\[
\boxed{
\hat{\beta}_0 \approx 1.41,
\quad
\hat{\beta}_1 \approx 0.39,
\quad
\hat{\beta}_2 \approx 2.16.
}
\]



\paragraph{5.1 SSE, SSR, SST.}

La \emph{Suma Total de Cuadrados} (SST) es
\[
\text{SST}
= \sum_{i=1}^5 (\,Y_i - \bar{Y}\,)^2.
\]
En este caso, la SST resulta $45.20$ (véase el modelo simple).  

El valor ajustado es
\[
\hat{Y}_i
= 1.41 + 0.39\,X_{1i} + 2.16\,X_{2i}.
\]

\[
\begin{array}{|c|c|c|c|c|c|c|}
\hline
\text{Ind} & X_1 & X_2 &Y 
& \hat{Y}_i = 1.41 + 0.39\,X_{1i} + 2.16\,X_{2i}.
& e_i= Y-\hat{Y}_i
& e_i^2
\\\hline
1 & 4 & 1 &5 
& 1.41 +0.39\cdot4 + 2.16\cdot 1 =5.13
& 5-5.13=-0.13
& 0.0169
\\
2 & 3 &2 &7
& 1.41 +0.39\cdot3 + 2.16\cdot2=6.9
& 7-6.9=0.1
& 0.01
\\
3 & 5 &3 &10
& 1.41 +0.39\cdot5 + 2.16\cdot3=9.84
& 10-9.84=0.16
& 0.0256
\\
4 & 7 &4 &13
& 1.41 +0.39\cdot4 + 2.16\cdot4=12.78
& 13-12.78=-0.22
& 0.0484
\\
5 & 6 &4 &12
& 1.41 +0.39\cdot4 + 2.16\cdot4=12.39
& 12-12.39=-0.39
& 0.1521
\\
\hline
& & & & & \sum e_i^2= & 0.253
\\\hline
\end{array}
\]

Los residuos $e_i=Y_i-\hat{Y}_i$ dan la \emph{Suma de Cuadrados de Residuos} (SSR). Tras los cálculos numéricos (ver tabla de $e_i^2$), se obtiene
\[
\text{SSR} \approx 0.253.
\]
Entonces
\[
\text{SSE}
= \text{SST} - \text{SSR}
= 45.20 - 0.253
= 44.947.
\]
\[
R^2
= 1-\frac{\text{SSR}}{\text{SST}}
= 1 - \frac{0.253}{45.20}
\approx 0.9944.
\]

\paragraph{5.2 $R_{\text{ajustado}}^2$.}

Con $n=5$ y $k=2$ variables explicativas (sin contar la constante), la fórmula es:
\[
R_{\text{ajustado}}^2
= 1
- 
\frac{\text{SSR}/(n-k-1)}{\text{SST}/(n-1)},
\]
es decir
\[
= 1
- 
\frac{0.253/(5-2-1)}{45.20/(5-1)}
= 1
- 
\frac{0.253/2}{45.20/4}.
\]
\[
0.253/2 = 0.1265,
\quad
45.20/4=11.30,
\quad
\frac{0.1265}{11.30}\approx 0.0112,
\]
\[
R_{\text{ajustado}}^2
= 1-0.0112
= 0.9888.
\]
\[
\boxed{
R^2 \approx 0.9944,
\quad
R_{\text{ajustado}}^2 \approx 0.9888.
}
\]

\paragraph{4.1 C\'alculos en R.} Usamos \texttt{lm()}:

\begin{lstlisting}
# Datos
X1 <- c(4,3,5,7,6)
X2 <- c(1,2,3,4,4)
Y  <- c(5,7,10,13,12)
datos <- data.frame(X1, X2, Y)

# Modelo simple:
mod_simple <- lm(Y ~ X1, data=datos)
summary(mod_simple)
# (Muestra coeficientes, R^2, etc.)

# Modelo multiple:
mod_multiple <- lm(Y ~ X1 + X2, data=datos)
summary(mod_multiple)
# (Nuevos coeficientes, R^2 mayor, etc.)
\end{lstlisting}

\item ¿De que valor es el sesgo por variable omitida en este ejercicio? ¿Como es la correlación entre $X_1$ y $X_2$
\\
\textbf{Solución}
\\
Vemos que el beta en el modelo simple es de $\tilde{\hat{\beta_1}}=1.90$. Y en el modelo múltiple es $\hat{\beta_1}=0.39$, entonces el sesgo esta dado por:
$$\tilde{\hat{\beta_1}}=\hat{\beta_1}+\hat\beta_2\cdot \delta \implies 1.9 = 0.39 + \hat\beta_2\cdot \delta$$
$$\implies \hat\beta_2\cdot \delta = 1.9-0.39=1.51$$
El sesgo es positivo de valor 1.51, por lo que, como $\beta_2$ es positivo, entonces $\delta$ es positivo, y esto indica que la correlación entre $X_1$ y $X_2$ es positiva.

\end{enumerate}




\item En una empresa se realizó un curso de un software importante para el desarrollo del trabajo, siendo este obligatorio con selección aleatoria para algunos trabajadores, y tambien los no escogidos aleatoriamente podían inscribirse de forma voluntaria.
Sea $Y$ el salario, $D\in\{0,1\}$ el tratamiento (tomó el curso) y $R\in\{0,1\}$ la asignación aleatoria.
Se observan 8 individuos (4 tratados, 4 no tratados). Entre los tratados, 3 son aleatorios ($R=1$); entre los no tratados, 2 son aleatorios.

\begin{center}
\begin{tabular}{|c| c |c |c|}
\hline
id & $Y$ & $D$ & $R$\\
\hline
1 & 950 & 1 & 1\\
2 & 900 & 1 & 1\\
3 & 880 & 1 & 1\\
4 & 820 & 1 & 0\\
5 & 780 & 0 & 1\\
6 & 740 & 0 & 1\\
7 & 700 & 0 & 0\\
8 & 650 & 0 & 0\\
\hline
\end{tabular}
\end{center}

\begin{enumerate}
\item Calcule la diferencia ingenua $\hat\beta_{\text{ing}}=\overline{Y}_{D=1}-\overline{Y}_{D=0}$.
\\
\textbf{Solución}
\\
El beta ingenuo (sesgado, antes de considerar aleatoriedad) esta dado por:
$$\beta_{ing}=E[Y|D=1]-E[Y|D=0]$$
$$E[Y|D=1] = \frac{950+900+880+820}{4}=887.5$$
$$E[Y|D=0]=\frac{780+740+700+650}{4}=717.5$$
$$\beta_{ing}=887.5-717.5=170$$

\item En la submuestra $R=1$, calcule $\hat\beta_{R} = \overline{Y}_{D=1,R=1}-\overline{Y}_{D=0,R=1}$.
\\
\textbf{Solución}
\\
Ahora considerando la aleatoriedad, consideraremos solo los casos que fueron escogidos aleatoriamente y calcularemos el $\beta$.
$$\hat\beta = E[{Y}|{D=1,R=1}]-E[{Y}|{D=0,R=1}]$$
$$ E[{Y}|{D=1,R=1}] = \frac{950+900+880}{3}=910$$
$$ E[{Y}|{D=0,R=1}] = \frac{780+740}{2}=760$$
$$\hat{\beta}= 910-760=150$$
\item Interprete $\hat\beta$ como estimador cercano al $ATE$ (y al $ATT$ con efectos homogéneos). Calcule el sesgo de selección en este caso.
\\
\textbf{Solución}
\\
Como el calculo de $\hat{\beta}$ es sobre aleatorios, podemos decir que $\hat{\beta}=ATE=ATT$, y si comparamos con el valor ingenuo, podemos encontrar el valor del sesgo:
$$Sesgo=\hat{\beta_{ing}}-\hat{\beta}=170-150=20$$
El sesgo era igual a 20, y estaba sobreestimando el valor real de $\beta$.
\item Discuta por qué $\hat\beta_{\text{ing}}$ puede estar sesgado cuando $D$ no es aleatorio.
\\
\textbf{Solución}
\\
La decisión de participar de un curso, cuando es de forma voluntaria, depende de varios factores, como inteligencia del individuo, motivación, entre otros... Y estas características podrían influir tambien en tener un salario mayor.
\end{enumerate}

\item Queremos realizar la siguiente regresión de $Y=\beta_0+\beta_1X+\epsilon$, pero solo podemos tener una versión aproximada de $Y^*=Y+c$, con $c$ una constante desconocida, y $X^*=X+u$, siendo $u$ error clásico ($cov(X,u)=0, var(u)=\sigma^2_u$), por lo que sólo podemos tener una versión imprecisa de la regresión, $Y^* = \tilde\beta_0+\tilde\beta_1X^*+v$. Desmuestra que 
$$\hat{\tilde\beta}_1=\hat{\beta_1}\frac{Var(X)}{Var(X)+\sigma^2_u}$$
\\
\textbf{Solución}
\\
necesitamos encontrar 
$$\tilde{\hat{\beta}}_1=\frac{Cov(X^*,Y^*)}{Var(X^*)}$$
sabemos que $Y^*=Y+c$, con $c$ una constante desconocida, y $X^*=X+u$, siendo $u$ error clásico. Entonces
$$\tilde{\hat{\beta}}_1=\frac{Cov(X+u,Y+c)}{Var(X+u)}$$
$$\tilde{\hat{\beta}}_1=\frac{Cov(X,Y)+Cov(X,c)+Cov(Y,u)+ Cov(c,u))}{Var(X)+Var(u)}$$
Sabemos que la varianza de una constante con cualquier valor es 0, y además como $u$ es error clásico, no correlaciona con $Y$, por lo que queda
$$\tilde{\hat{\beta}}_1=\frac{Cov(X,Y))}{Var(X)+\sigma^2_u}$$
Luego sabemos que $Y=\beta_0+\beta_1X+\epsilon$, por lo que
$$\tilde{\hat{\beta}}_1=\frac{Cov(X,\beta_0+\beta_1X+\epsilon))}{Var(X)+\sigma^2_u}$$
distribuyendo:
$$\tilde{\hat{\beta}}_1=\frac{Cov(X,\beta_0)+Cov(X,\beta_1X)+Cov(X,\epsilon))}{Var(X)+\sigma^2_u}$$
$Cov(X,\beta_0)=0$ porque $\beta_0$ es constante y $Cov(X,\epsilon)=0$, por el supuesto de $E[\epsilon|X]=0$, entonces:
$$\tilde{\hat{\beta}}_1=\frac{Cov(X,\beta_1X)}{Var(X)+\sigma^2_u}=\beta_1\frac{Var(X)}{Var(X)+\sigma^2_u}$$
\end{enumerate}

\section{Formulario}


\begin{itemize}

\item  Estimadores MCO para RLS.
{
$$\hat{\beta}_0 = \bar{y} - \hat{\beta}_1 \bar{x} $$
$$\hat{\beta}_1 = \frac{\sum_{i=1}^{n}(x_i - \bar{x})(y_i-\bar{y})}{\sum_{i=1}^{n}(x_i - \bar{x})^2}=\frac{Cov(x,y)}{Var(x)}$$}


\item  Estimadores MCO para RLM.
{
$${\hat{\beta}} = (X^TX)^{-1} X^TY $$}


    \item \[ R^2  = 1 - \frac{SSR}{SST} \] 
\item  $$SST = \sum_{i=1}^{n} (y_i - \bar{y})^2 \quad \quad SSR = \sum_{i=1}^{n}  \hat{u}_i^2 $$
    \item \[ \bar{R}^2 =  1 - \frac{(n-1)SSR}{(n-k-1)SST} \]

    \item  $$\hat{\sigma}^2 = \frac{\sum_{i=1}^{n} \hat{u}_i^2 }{n-k-1}= \frac{SSR}{n - k - 1}$$
    \item \[
\text{Var}(\hat{\beta}|X) = \hat{\sigma}^2(X^{\top}X)^{-1}
\]
    \end{itemize}



\end{document}